\IfFileExists{stacks-project.cls}{%
\documentclass{stacks-project}
}{%
\documentclass{amsart}
}

% For dealing with references we use the comment environment
\usepackage{verbatim}
\newenvironment{reference}{\comment}{\endcomment}
%\newenvironment{reference}{}{}
\newenvironment{slogan}{\comment}{\endcomment}
\newenvironment{history}{\comment}{\endcomment}

% For commutative diagrams we use Xy-pic
\usepackage[all]{xy}

% We use 2cell for 2-commutative diagrams.
\xyoption{2cell}
\UseAllTwocells

% We use multicol for the list of chapters between chapters
\usepackage{multicol}

% This is generally recommended for better output
\usepackage{lmodern}
\usepackage[T1]{fontenc}

% For cross-file-references
\usepackage{xr-hyper}

% Package for hypertext links:
\usepackage{hyperref}

% For any local file, say "hello.tex" you want to link to please
% use \externaldocument[hello-]{hello}
\externaldocument[introduction-]{introduction}
\externaldocument[conventions-]{conventions}
\externaldocument[sets-]{sets}
\externaldocument[categories-]{categories}
\externaldocument[topology-]{topology}
\externaldocument[sheaves-]{sheaves}
\externaldocument[sites-]{sites}
\externaldocument[stacks-]{stacks}
\externaldocument[fields-]{fields}
\externaldocument[algebra-]{algebra}
\externaldocument[brauer-]{brauer}
\externaldocument[homology-]{homology}
\externaldocument[derived-]{derived}
\externaldocument[simplicial-]{simplicial}
\externaldocument[more-algebra-]{more-algebra}
\externaldocument[smoothing-]{smoothing}
\externaldocument[modules-]{modules}
\externaldocument[sites-modules-]{sites-modules}
\externaldocument[injectives-]{injectives}
\externaldocument[cohomology-]{cohomology}
\externaldocument[sites-cohomology-]{sites-cohomology}
\externaldocument[dga-]{dga}
\externaldocument[dpa-]{dpa}
\externaldocument[sdga-]{sdga}
\externaldocument[hypercovering-]{hypercovering}
\externaldocument[schemes-]{schemes}
\externaldocument[constructions-]{constructions}
\externaldocument[properties-]{properties}
\externaldocument[morphisms-]{morphisms}
\externaldocument[coherent-]{coherent}
\externaldocument[divisors-]{divisors}
\externaldocument[limits-]{limits}
\externaldocument[varieties-]{varieties}
\externaldocument[topologies-]{topologies}
\externaldocument[descent-]{descent}
\externaldocument[perfect-]{perfect}
\externaldocument[more-morphisms-]{more-morphisms}
\externaldocument[flat-]{flat}
\externaldocument[groupoids-]{groupoids}
\externaldocument[more-groupoids-]{more-groupoids}
\externaldocument[etale-]{etale}
\externaldocument[chow-]{chow}
\externaldocument[intersection-]{intersection}
\externaldocument[pic-]{pic}
\externaldocument[weil-]{weil}
\externaldocument[adequate-]{adequate}
\externaldocument[dualizing-]{dualizing}
\externaldocument[duality-]{duality}
\externaldocument[discriminant-]{discriminant}
\externaldocument[derham-]{derham}
\externaldocument[local-cohomology-]{local-cohomology}
\externaldocument[algebraization-]{algebraization}
\externaldocument[curves-]{curves}
\externaldocument[resolve-]{resolve}
\externaldocument[models-]{models}
\externaldocument[functors-]{functors}
\externaldocument[equiv-]{equiv}
\externaldocument[pione-]{pione}
\externaldocument[etale-cohomology-]{etale-cohomology}
\externaldocument[proetale-]{proetale}
\externaldocument[relative-cycles-]{relative-cycles}
\externaldocument[more-etale-]{more-etale}
\externaldocument[trace-]{trace}
\externaldocument[crystalline-]{crystalline}
\externaldocument[spaces-]{spaces}
\externaldocument[spaces-properties-]{spaces-properties}
\externaldocument[spaces-morphisms-]{spaces-morphisms}
\externaldocument[decent-spaces-]{decent-spaces}
\externaldocument[spaces-cohomology-]{spaces-cohomology}
\externaldocument[spaces-limits-]{spaces-limits}
\externaldocument[spaces-divisors-]{spaces-divisors}
\externaldocument[spaces-over-fields-]{spaces-over-fields}
\externaldocument[spaces-topologies-]{spaces-topologies}
\externaldocument[spaces-descent-]{spaces-descent}
\externaldocument[spaces-perfect-]{spaces-perfect}
\externaldocument[spaces-more-morphisms-]{spaces-more-morphisms}
\externaldocument[spaces-flat-]{spaces-flat}
\externaldocument[spaces-groupoids-]{spaces-groupoids}
\externaldocument[spaces-more-groupoids-]{spaces-more-groupoids}
\externaldocument[bootstrap-]{bootstrap}
\externaldocument[spaces-pushouts-]{spaces-pushouts}
\externaldocument[spaces-chow-]{spaces-chow}
\externaldocument[groupoids-quotients-]{groupoids-quotients}
\externaldocument[spaces-more-cohomology-]{spaces-more-cohomology}
\externaldocument[spaces-simplicial-]{spaces-simplicial}
\externaldocument[spaces-duality-]{spaces-duality}
\externaldocument[formal-spaces-]{formal-spaces}
\externaldocument[restricted-]{restricted}
\externaldocument[spaces-resolve-]{spaces-resolve}
\externaldocument[formal-defos-]{formal-defos}
\externaldocument[defos-]{defos}
\externaldocument[cotangent-]{cotangent}
\externaldocument[examples-defos-]{examples-defos}
\externaldocument[algebraic-]{algebraic}
\externaldocument[examples-stacks-]{examples-stacks}
\externaldocument[stacks-sheaves-]{stacks-sheaves}
\externaldocument[criteria-]{criteria}
\externaldocument[artin-]{artin}
\externaldocument[quot-]{quot}
\externaldocument[stacks-properties-]{stacks-properties}
\externaldocument[stacks-morphisms-]{stacks-morphisms}
\externaldocument[stacks-limits-]{stacks-limits}
\externaldocument[stacks-cohomology-]{stacks-cohomology}
\externaldocument[stacks-perfect-]{stacks-perfect}
\externaldocument[stacks-introduction-]{stacks-introduction}
\externaldocument[stacks-more-morphisms-]{stacks-more-morphisms}
\externaldocument[stacks-geometry-]{stacks-geometry}
\externaldocument[moduli-]{moduli}
\externaldocument[moduli-curves-]{moduli-curves}
\externaldocument[examples-]{examples}
\externaldocument[exercises-]{exercises}
\externaldocument[guide-]{guide}
\externaldocument[desirables-]{desirables}
\externaldocument[coding-]{coding}
\externaldocument[obsolete-]{obsolete}
\externaldocument[fdl-]{fdl}
\externaldocument[index-]{index}

% Theorem environments.
%
\theoremstyle{plain}
\newtheorem{theorem}[subsection]{Theorem}
\newtheorem{proposition}[subsection]{Proposition}
\newtheorem{lemma}[subsection]{Lemma}

\theoremstyle{definition}
\newtheorem{definition}[subsection]{Definition}
\newtheorem{example}[subsection]{Example}
\newtheorem{exercise}[subsection]{Exercise}
\newtheorem{situation}[subsection]{Situation}

\theoremstyle{remark}
\newtheorem{remark}[subsection]{Remark}
\newtheorem{remarks}[subsection]{Remarks}

\numberwithin{equation}{subsection}

% Macros
%
\def\lim{\mathop{\mathrm{lim}}\nolimits}
\def\colim{\mathop{\mathrm{colim}}\nolimits}
\def\Spec{\mathop{\mathrm{Spec}}}
\def\Hom{\mathop{\mathrm{Hom}}\nolimits}
\def\Ext{\mathop{\mathrm{Ext}}\nolimits}
\def\SheafHom{\mathop{\mathcal{H}\!\mathit{om}}\nolimits}
\def\SheafExt{\mathop{\mathcal{E}\!\mathit{xt}}\nolimits}
\def\Sch{\mathit{Sch}}
\def\Mor{\mathop{\mathrm{Mor}}\nolimits}
\def\Ob{\mathop{\mathrm{Ob}}\nolimits}
\def\Sh{\mathop{\mathit{Sh}}\nolimits}
\def\NL{\mathop{N\!L}\nolimits}
\def\CH{\mathop{\mathrm{CH}}\nolimits}
\def\proetale{{pro\text{-}\acute{e}tale}}
\def\etale{{\acute{e}tale}}
\def\QCoh{\mathit{QCoh}}
\def\Ker{\mathop{\mathrm{Ker}}}
\def\Im{\mathop{\mathrm{Im}}}
\def\Coker{\mathop{\mathrm{Coker}}}
\def\Coim{\mathop{\mathrm{Coim}}}

% Boxtimes
%
\DeclareMathSymbol{\boxtimes}{\mathbin}{AMSa}{"02}

%
% Macros for moduli stacks/spaces
%
\def\QCohstack{\mathcal{QC}\!\mathit{oh}}
\def\Cohstack{\mathcal{C}\!\mathit{oh}}
\def\Spacesstack{\mathcal{S}\!\mathit{paces}}
\def\Quotfunctor{\mathrm{Quot}}
\def\Hilbfunctor{\mathrm{Hilb}}
\def\Curvesstack{\mathcal{C}\!\mathit{urves}}
\def\Polarizedstack{\mathcal{P}\!\mathit{olarized}}
\def\Complexesstack{\mathcal{C}\!\mathit{omplexes}}
% \Pic is the operator that assigns to X its picard group, usage \Pic(X)
% \Picardstack_{X/B} denotes the Picard stack of X over B
% \Picardfunctor_{X/B} denotes the Picard functor of X over B
\def\Pic{\mathop{\mathrm{Pic}}\nolimits}
\def\Picardstack{\mathcal{P}\!\mathit{ic}}
\def\Picardfunctor{\mathrm{Pic}}
\def\Deformationcategory{\mathcal{D}\!\mathit{ef}}
\usepackage{cleveref}
\crefname{originalreference}{original reference}{original references}
\Crefname{originalreference}{Original reference}{Original references}
\makeatletter
\newlabel{cohomology-lemma-dual-perfect-complex}{{08DQ}{source}{}{sourcepointer-1}{}}
\newlabel{cohomology-lemma-dual-perfect-complex@cref}{{[originalreference][0][]08DQ}{[1][0][]source}{}{}{}}
\makeatother
\crefname{originalreference}{original reference}{original references}
\Crefname{originalreference}{Original reference}{Original references}
\makeatletter
\newlabel{dga-lemma-upgrade-tensor-with-complex-derived}{{0CS6}{source}{}{sourcepointer-2}{}}
\newlabel{dga-lemma-upgrade-tensor-with-complex-derived@cref}{{[originalreference][0][]0CS6}{[1][0][]source}{}{}{}}
\makeatother
\crefname{originallemma}{original lemma}{original lemmas}
\Crefname{originallemma}{Original lemma}{Original lemmas}
\makeatletter
\newlabel{lemma-perfect-for-R}{{0FYA}{source}{}{sourcepointer-3}{}}
\newlabel{lemma-perfect-for-R@cref}{{[originallemma][0][]0FYA}{[1][0][]source}{}{}{}}
\makeatother
\crefname{originalreference}{original reference}{original references}
\Crefname{originalreference}{Original reference}{Original references}
\makeatletter
\newlabel{perfect-lemma-Pn-module-category}{{0BQU}{source}{}{sourcepointer-4}{}}
\newlabel{perfect-lemma-Pn-module-category@cref}{{[originalreference][0][]0BQU}{[1][0][]source}{}{}{}}
\makeatother
\crefname{originalreference}{original reference}{original references}
\Crefname{originalreference}{Original reference}{Original references}
\makeatletter
\newlabel{perfect-lemma-direct-image-coherent}{{08E2}{source}{}{sourcepointer-5}{}}
\newlabel{perfect-lemma-direct-image-coherent@cref}{{[originalreference][0][]08E2}{[1][0][]source}{}{}{}}
\makeatother
\crefname{originalreference}{original reference}{original references}
\Crefname{originalreference}{Original reference}{Original references}
\makeatletter
\newlabel{perfect-lemma-perfect-on-regular}{{0FDC}{source}{}{sourcepointer-6}{}}
\newlabel{perfect-lemma-perfect-on-regular@cref}{{[originalreference][0][]0FDC}{[1][0][]source}{}{}{}}
\makeatother
\crefname{originalreference}{original reference}{original references}
\Crefname{originalreference}{Original reference}{Original references}
\makeatletter
\newlabel{perfect-proposition-compact-is-perfect}{{09M1}{source}{}{sourcepointer-7}{}}
\newlabel{perfect-proposition-compact-is-perfect@cref}{{[originalreference][0][]09M1}{[1][0][]source}{}{}{}}
\makeatother
\title{Boundedness criterion on projective space via the Beilinson fan}
\author{The Stacks Project Authors}
\date{stacks-project:equiv.tex:local-snapshot-20261012; GNU FDL 1.2 or later}
\begin{document}
\maketitle
\noindent\textit{Allowed standard prerequisites (editorial).}\begin{itemize}
\item Perfect duality. cohomology-lemma-dual-perfect-complex (cross-chapter, named). (2)-(3) equivalence and the Ext reidentification.
\item Endomorphism ring equivalence. perfect-lemma-Pn-module-category (cross-chapter, named). The Beilinson fan P gives an equivalence D(R) -> D\_QCoh(O\_\{P\textasciicircum{}n\}).
\item Endomorphism ring compactness. lemma-perfect-for-R L26459 (carded this batch). Condition (2) of that lemma matches the displayed Ext computation.
\end{itemize}\medskip
\ifdefined\hypertarget\hypertarget{sourcepointer-1}{}\fi
\noindent\textit{Source reference.} Original reference 08DQ: focused-seeds/geometry/S9-eqv-coherent-on-projective-space/ref-tags.txt, source lines 5--5. Navigation pointer only: the counted proof cites this statement by \textbackslash{}ref as a standing input; the deduction the counted proof has to make is carried out inside the counted span, and no part of that deduction is delegated to this pointer.\par\medskip
\ifdefined\hypertarget\hypertarget{sourcepointer-2}{}\fi
\noindent\textit{Source reference.} Original reference 0CS6: focused-seeds/geometry/S9-eqv-coherent-on-projective-space/ref-tags.txt, source lines 6--6. Navigation pointer only: the counted proof cites this statement by \textbackslash{}ref as a standing input; the deduction the counted proof has to make is carried out inside the counted span, and no part of that deduction is delegated to this pointer.\par\medskip
\ifdefined\hypertarget\hypertarget{sourcepointer-3}{}\fi
\noindent\textit{Source reference.} Original lemma 0FYA: sources/stacks-equiv/pinned-original-source.tex, source lines 676--676. Navigation pointer only: the counted proof cites this statement by \textbackslash{}ref as a standing input; the deduction the counted proof has to make is carried out inside the counted span, and no part of that deduction is delegated to this pointer.\par\medskip
\ifdefined\hypertarget\hypertarget{sourcepointer-4}{}\fi
\noindent\textit{Source reference.} Original reference 0BQU: focused-seeds/geometry/S9-eqv-coherent-on-projective-space/ref-tags.txt, source lines 7--7. Navigation pointer only: the counted proof cites this statement by \textbackslash{}ref as a standing input; the deduction the counted proof has to make is carried out inside the counted span, and no part of that deduction is delegated to this pointer.\par\medskip
\ifdefined\hypertarget\hypertarget{sourcepointer-5}{}\fi
\noindent\textit{Source reference.} Original reference 08E2: focused-seeds/geometry/S9-eqv-coherent-on-projective-space/ref-tags.txt, source lines 8--8. Navigation pointer only: the counted proof cites this statement by \textbackslash{}ref as a standing input; the deduction the counted proof has to make is carried out inside the counted span, and no part of that deduction is delegated to this pointer.\par\medskip
\ifdefined\hypertarget\hypertarget{sourcepointer-6}{}\fi
\noindent\textit{Source reference.} Original reference 0FDC: focused-seeds/geometry/S9-eqv-coherent-on-projective-space/ref-tags.txt, source lines 9--9. Navigation pointer only: the counted proof cites this statement by \textbackslash{}ref as a standing input; the deduction the counted proof has to make is carried out inside the counted span, and no part of that deduction is delegated to this pointer.\par\medskip
\ifdefined\hypertarget\hypertarget{sourcepointer-7}{}\fi
\noindent\textit{Source reference.} Original reference 09M1: focused-seeds/geometry/S9-eqv-coherent-on-projective-space/ref-tags.txt, source lines 10--10. Navigation pointer only: the counted proof cites this statement by \textbackslash{}ref as a standing input; the deduction the counted proof has to make is carried out inside the counted span, and no part of that deduction is delegated to this pointer.\par\medskip
\begin{lemma}
\label{lemma-coherent-on-projective-space}
Let $k$ be a field. Let $n \geq 0$. Let
$K \in D_\QCoh(\mathcal{O}_{\mathbf{P}^n_k})$.
The following are equivalent
\begin{enumerate}
\item $K$ is in $D^b_{\textit{Coh}}(\mathcal{O}_{\mathbf{P}^n_k})$,
\item $\sum_{i \in \mathbf{Z}}
\dim_k H^i(\mathbf{P}^n_k, E \otimes^\mathbf{L} K) < \infty$
for each perfect object $E$ of
$D(\mathcal{O}_{\mathbf{P}^n_k})$,
\item $\sum_{i \in \mathbf{Z}}
\dim_k \Ext^i_{\mathbf{P}^n_k}(E, K) < \infty$
for each perfect object $E$ of $D(\mathcal{O}_{\mathbf{P}^n_k})$,
\item $\sum_{i \in \mathbf{Z}} \dim_k H^i(\mathbf{P}^n_k,
K \otimes^\mathbf{L} \mathcal{O}_{\mathbf{P}^n_k}(d)) < \infty$
for $d = 0, 1, \ldots, n$.
\end{enumerate}
\end{lemma}

\begin{proof}
Parts (2) and (3) are equivalent by
Cohomology, Lemma \ref{cohomology-lemma-dual-perfect-complex}.
If (1) is true, then for $E$ perfect the derived tensor product
$E \otimes^\mathbf{L} K$ is in
$D^b_{\textit{Coh}}(\mathcal{O}_{\mathbf{P}^n_k})$
and we see that (2) holds by 
Derived Categories of Schemes, Lemma \ref{perfect-lemma-direct-image-coherent}.
It is clear that (2) implies (4) as $\mathcal{O}_{\mathbf{P}^n_k}(d)$
can be viewed
as a perfect object of the derived category of $\mathbf{P}^n_k$.
Thus it suffices to prove that (4) implies (1).

\medskip\noindent
Assume (4). Let $R$ be as in Lemma \ref{lemma-perfect-for-R}.
Let $P = \bigoplus_{d = 0, \ldots, n} \mathcal{O}_{\mathbf{P}^n_k}(-d)$.
Recall that $R = \text{End}_{\mathbf{P}^n_k}(P)$ whereas all other
self-Exts of $P$ are zero and that $P$ determines an equivalence
$- \otimes^\mathbf{L} P : D(R) \to D_\QCoh(\mathcal{O}_{\mathbf{P}^n_k})$
by Derived Categories of Schemes, Lemma \ref{perfect-lemma-Pn-module-category}.
Say $K$ corresponds to $L$ in $D(R)$. Then
\begin{align*}
H^i(L)
& =
\Ext^i_{D(R)}(R, L) \\
& =
\Ext^i_{\mathbf{P}^n_k}(P, K) \\
& =
H^i(\mathbf{P}^n_k, K \otimes P^\vee) \\
& =
\bigoplus\nolimits_{d = 0, \ldots, n}
H^i(\mathbf{P}^n_k, K \otimes \mathcal{O}(d))
\end{align*}
by Differential Graded Algebra, Lemma
\ref{dga-lemma-upgrade-tensor-with-complex-derived}
(and the fact that $- \otimes^\mathbf{L} P$ is an equivalence)
and Cohomology, Lemma \ref{cohomology-lemma-dual-perfect-complex}.
Thus our assumption (4) implies that $L$ satisfies condition (2) of
Lemma \ref{lemma-perfect-for-R} and hence is a compact object of $D(R)$.
Therefore $K$ is a compact object of
$D_\QCoh(\mathcal{O}_{\mathbf{P}^n_k})$.
Thus $K$ is perfect by
Derived Categories of Schemes, Proposition
\ref{perfect-proposition-compact-is-perfect}.
Since $D_{perf}(\mathcal{O}_{\mathbf{P}^n_k}) =
D^b_{\textit{Coh}}(\mathcal{O}_{\mathbf{P}^n_k})$
by
Derived Categories of Schemes, Lemma \ref{perfect-lemma-perfect-on-regular}
we conclude (1) holds.
\end{proof}

\end{document}
